\paragraph{What scRareRefine does and does not solve.}
scRareRefine addresses a specific and well-defined sub-problem: recovering known rare cell types from scANVI-based annotation when the rare class is geometrically separable in latent space but the softmax classifier is miscalibrated due to extreme label imbalance. Within this scope, the results are strong and consistent. Outside this scope, the framework correctly abstains: for low-separability cell types or already-reliable baselines, it produces near-zero modification by design.

The method is \emph{not} a general-purpose annotation tool or a replacement for scANVI. It does not address: (i) \emph{de novo} discovery of entirely unknown rare populations, which requires unsupervised anomaly detection \citep{xu2024scCAD,fang2021gapclust}; (ii) cell types that are not geometrically separable in scANVI's latent space, for which improved embeddings or retraining would be needed; or (iii) scenarios where marker-gene expression is highly stochastic or batch-dependent, undermining the biological verification stage.

\paragraph{Relationship to logistic regression.}
An important finding is that logistic regression (LR), applied to normalized expression features rather than latent embeddings, outperforms scRareRefine on low-separability cell types such as $\varepsilon$-cells. This is consistent with the separability ratio framework: when the rare class is not well-separated in latent space, expression-space classifiers may exploit gene-level features that are obscured by the batch-correction step in scANVI. The complementarity between scRareRefine and LR suggests a practical ensemble strategy: use $\sepratio$ to decide which method to deploy, or combine both predictions.

\paragraph{Scope of the evaluation.}
The current evaluation spans seven datasets and nine rare cell types. While the consistency of the separability-ratio pattern across immune, pancreatic, liver, spleen, kidney, and PBMC data is encouraging, the number of cell-type-level data points remains limited (Spearman $\rho = 0.81$, $n = 9$). Conclusions about the generality of the $\sepratio > 1.3$ threshold should be treated as preliminary until validated on a larger benchmark covering more tissue types, species, and sequencing platforms. We note that the threshold values (1.1 and 1.3) were not tuned on the evaluation data; they emerged from prior exploratory experiments and are reported here without post-hoc adjustment.

\paragraph{Comparison scope.}
The current evaluation does not include scBalance \citep{xu2023scbalance}, sc-SynO \citep{mulvey2021scsyno}, or foundation model embeddings (scGPT \citep{cui2024scgpt}) as baselines. A systematic comparison against imbalance-aware retraining methods and large foundation model representations is a priority for future work. We anticipate that scBalance and sc-SynO would show competitive performance at larger annotation budgets but may not be able to match the prototype rescue performance at $n_\text{rare} = 5$, where synthetic augmentation is most unreliable.

\paragraph{Sensitivity to hyperparameters.}
scRareRefine involves several hyperparameters: the margin percentile threshold ($q_{25}$), the marker gene count ($k = 25$), the false-rescue rate tolerance ($\text{FRR} \leq 0.001$), and the separability thresholds (1.1, 1.3). A detailed sensitivity analysis is provided in \cref{sec:appendix}. In brief, results are relatively insensitive to moderate changes in these parameters; the most impactful choice is the separability threshold, with $\sepratio = 1.3$ providing a good balance between rescue activation and abstention.

\paragraph{Applicability beyond scANVI.}
Although all experiments use scANVI embeddings, the scRareRefine framework is applicable to any method that produces a fixed-dimensional latent embedding alongside a softmax cell-type prediction. Candidate upstream models include scArches \citep{lotfollahi2022mapping}, fine-tuned scGPT \citep{cui2024scgpt}, and CellPLM. The separability ratio would need to be recomputed for each embedding, but the pipeline logic is unchanged. Evaluating scRareRefine on foundation model embeddings is a natural extension of this work.
